\section{\label{sec:lattsmear} 格点上的混合}

在详细考察 sOPE 之前，我们先演示梯度流如何消除格点上的幂次发散混合。我们以标量场论中的
twist-2 算符为例；它们对称且无迹，由下式给出：
\begin{equation}
T_{\mu_1\ldots\mu_n}(x) = \phi(x)\partial_{\mu_1}\ldots
\partial_{\mu_n}\phi(x) - \mathrm{traces}.
\end{equation}
在格点上，我们把偏导数替换为离散差分算符
\begin{equation}
T_{\mu_1\ldots\mu_n}^{\mathrm{latt}}(x) = \phi(x)\Delta_{\mu_1}\ldots
\Delta_{\mu_n}\phi(x) -
\mathrm{traces}.
\end{equation}
此时时空点 $x$ 是格点上的一个节点，$x_\mu =
an_\mu$，其中 $a$ 为格距，$n_\mu$ 是四分量格点矢量。离散差分算符可以加以改进以消除离散化效应，
而最简单的此类对称算符对标量场的作用为
\begin{equation}
\Delta_\mu \phi(x) = \frac{1}{2a}\left(\phi(x+\widehat{\mu}) -
\phi(x-\widehat{\mu})\right),
\end{equation}
其中 $\hat{\mu}$ 是第 $\mu^{\mathrm{th}}$ 方向的单位矢量。

格点调节子破坏转动对称性，从而在质量维数不同的 twist-2 算符之间诱导混合。由量纲分析，
混合系数以格距的负幂次标度，并在连续极限下发散。这就是格点上的幂次发散混合问题。例如，
简单的双线性算符 $T^{\mathrm{latt}}(x) = \phi(x) \phi(x)$ 以 $1/a^2$
标度的系数与算符 $T_{\mu\nu}^{\mathrm{latt}}(x)$ 混合。更一般地，
$T^{\mathrm{latt}}(x)$ 与任何含偶数个导数的 twist-2 算符、即
$T_{\mu_1\ldots\mu_{2n}}^{\mathrm{latt}}(x)$ 之间的混合以 $1/a^{2n}$ 标度。

这些算符的涂抹对应物由
\begin{equation}
S_{\mu_1\ldots\mu_n}^{\mathrm{latt}}(x) =
\rho(\tau,x)\Delta_{\mu_1}\ldots
\Delta_{\mu_n}\rho(\tau,x) -
\mathrm{traces}.
\end{equation}
给出，不存在这一问题。
$S^{\mathrm{latt}}(x) = \rho(\tau,x) \rho(\tau,x)$ 与
$S_{\mu_1\ldots\mu_{2n}}^{\mathrm{latt}}(x)$ 之间的混合系数以
$1/\tau^{n}$ 标度，其中 $\tau$ 是以\emph{物理单位}计的流时间。只要保持量纲量
$\tau=\widetilde{\tau}a^2$ 固定（$\widetilde{\tau}$ 无量纲），
则随着格距 $a$ 减小，混合系数保持有限，因为在非零流时间处的矩阵元不可能包含任何额外的发散
\cite{Makino:2014sta,Luscher:2011bx}。

作为这一行为的简单示例，我们来考虑质量维数不同的两个 twist-2 算符的矩阵元：
\begin{equation}
M_{\mathrm{cont}} = \left\langle \Omega | \phi(0) \phi(0) \,\phi(0)
\partial_\mu\partial_\nu \phi(0) |\Omega \right\rangle.
\end{equation}
对无质量标量场，该矩阵元为零。然而在格点上，相应的矩阵元
\begin{equation}
M_{\mathrm{latt}} = \left\langle \Omega | \phi(0) \phi(0) \,\phi(0)
\Delta_\mu\Delta_\nu \phi(0) |\Omega \right\rangle,
\end{equation}
却不为零。

我们在图~\ref{fig:mmix} 的左图中给出对该矩阵元领头贡献的费曼图，它由
\begin{equation}
M^{(0)}_{\mathrm{latt}} =
\frac{1}{a^4}\int^{\pi/a}_{-\pi/a}\frac{\mathrm{d}^4(ak)}{(2\pi)^4}\
\frac{\widehat{k}_\mu\widehat{k}_\nu}{\big(\widehat{k}^2+m_0^2\big)^2 } ,
\end{equation}
给出，其中 $\widehat{k}_\mu = (2/a)\sin(ak_\mu/2)$。我们已引入裸质量
$m_0$ 以去除伪红外发散，但该积分是有限的，我们可以取无质量极限
$m_0 \rightarrow 0$ 来匹配连续的无质量理论。
\begin{figure}[htbp]
\begin{minipage}{0.49\textwidth}
\includegraphics[width=0.95\textwidth,keepaspectratio=true]{images/mmix}
\end{minipage}%
\begin{minipage}{0.49\textwidth}
\includegraphics[width=0.95\textwidth,keepaspectratio=true]{images/mmix_smear}
\end{minipage}
\caption{\label{fig:mmix}表示无涂抹场混合矩阵元 $M_{\mathrm{latt}}^{(0)}$ 与涂抹场矩阵元
$M_{\mathrm{latt}}^{(0)}(\tau)$ 领头贡献的费曼图。左图中的黑色实线代表未涂抹传播子，黑色实心方块
和菱形分别是算符 $\phi^2(0)$ 与 $\phi(0)\Delta_\mu\Delta_\nu\phi(0)$。右图中的蓝色虚线为涂抹传播子，
两个不同的蓝色圆斑（实心与空心）分别表示涂抹算符
$\rho^2(\tau,0)$ 与 $\rho(\tau,0)\Delta_\mu\Delta_\nu\rho(\tau,0)$。}
\end{figure}
我们按格距的幂次展开该积分，得到
\begin{align}
M^{(0)}_{\mathrm{latt}} = {} & -\frac{\delta_{\mu\nu}}{4}
\int^{\pi/a}_{-\pi/a}\frac{\mathrm{d}^4k}{(2\pi)^4}
\frac{k^2}{(k^2+m_0^2)^2 } +{\cal O}(a^2) \nonumber\\
= {} &
-\frac{\delta_{\mu\nu}}{48a^2}+{\cal
O}(a^0)
\end{align}
（无质量极限）。这一结果标志着幂次发散混合的出现：$1/a^2$ 项在连续极限下发散。尽管此计算仅到微扰论
领头阶，更高阶的贡献并不改变这种对格距的幂次发散依赖 \cite{Davoudi:2012ya}。

然而，如果我们用涂抹自由度计算这个矩阵元（见图~\ref{fig:mmix} 右图），则会发现
\begin{align}
M^{(0)}_{\mathrm{latt}}(\tau) = {} &\left\langle \Omega | \rho(\tau,0)
\rho(\tau,0) \,\rho(\tau,0)
\Delta_\mu\Delta_\nu \rho(\tau,0) |\Omega \right\rangle \nonumber\\
= {} &  -\frac{\delta_{\mu\nu}}{4}
\int^{\pi/a}_{-\pi/a}\frac{\mathrm{d}^4k}{(2\pi)^4}
\frac{k^2e^{-4\tau k^2}}{(k^2+m_0^2)^2 } +{\cal O}(a^2) \nonumber\\
= {} & \delta_{\mu\nu}
\frac{e^{-4\pi^2\tau/a^2}-1}{256\pi^2\tau}+{\cal
O}(a^0).
\end{align}
这里我们同样取了无质量极限。在连续极限下、并保持流时间 $\tau$ 以物理单位固定时，
该矩阵元趋于一个常数，这表明对涂抹自由度而言幂次发散混合受到了压低。
